\chapter{Metodologia}

Este capítulo descreve a arquitetura técnica e os procedimentos metodológicos adotados para a construção do sistema de detecção de \textit{fake news}. A pesquisa possui caráter aplicado e quantitativo, dividindo-se em três pilares principais de engenharia: (1) Coleta e processamento de dados do Bluesky via \textit{AT Protocol}; (2) Treinamento e modelagem da Inteligência Artificial em grafos; e (3) Desenvolvimento da aplicação web de ponta a ponta.

\section{Arquitetura Geral do Sistema}

O ecossistema do projeto foi desenhado para processar informações estruturadas em rede e exibi-las de forma intuitiva. O fluxo de dados inicia-se na extração de postagens e interações da rede social, que são transformadas em grafos e submetidas ao modelo preditivo. O resultado dessa inferência é então servido por uma API RESTful para consumo da interface do usuário.

\section{Engenharia de Dados e a Transição Metodológica}

O desenho da pesquisa passou por uma adaptação crucial em sua fase inicial, refletindo os desafios contemporâneos da coleta de dados em plataformas digitais fechadas.

\subsection{O Pivot do Projeto: Do Twitter ao Bluesky}

Inicialmente, a pesquisa visava utilizar o Twitter (atualmente X) como ambiente primário de análise, baseando os primeiros testes no \textit{benchmark} UPFD (\textit{User Preference-aware Fake News Detection}). No entanto, as recentes mudanças nas políticas de monetização da plataforma impuseram severas restrições comerciais de acesso à sua API, tornando inviável a extração de dados topológicos em larga escala por pesquisadores independentes.

Diante dessa barreira técnica e financeira, o projeto realizou um \textit{pivot} metodológico estratégico para a rede social Bluesky. A escolha foi motivada pela disponibilidade de um \textit{dataset} massivo de 21 GB contendo o histórico estrutural necessário, aliado ao fato de a plataforma possuir uma API aberta e gratuita baseada no \textit{AT Protocol}.

\subsection{O Bluesky como Microcosmo de Experimentação}

Embora o Bluesky possua uma base de usuários consideravelmente menor em comparação a redes consolidadas, essa característica provou-se uma vantagem metodológica. A rede funcionou como um excelente ambiente controlado de experimentação (um microcosmo topológico). Por ser uma rede em fase de adoção e com alta instabilidade na formação de laços sociais, o Bluesky simula perfeitamente o cenário de \textit{Cold Start} (Partida Fria). Nesse contexto, o algoritmo é forçado a classificar a propagação da informação sem depender de um histórico longo e consolidado dos usuários, testando os limites e a real eficácia das Redes Neurais de Grafos na detecção precoce de anomalias.

\subsection{Extração e Estruturação do Dataset Massivo}

O \textit{dataset} bruto de 21 GB armazenava dados de interação do Bluesky compactados em múltiplos formatos (\texttt{.tar.gz}, \texttt{.csv.gz} e \texttt{.jsonl.gz}). Para torná-lo operacional, implementou-se um módulo de extração seletiva (\texttt{extrator\_amostra.py}) que realiza a leitura em modo \textit{streaming}, sem descompactar o arquivo integralmente na memória.

O módulo opera em duas etapas: primeiramente, extrai a totalidade dos posts categorizados por feed temático do arquivo \texttt{feed\_posts.tar.gz}, que contém publicações organizadas em 11 categorias (\textit{News}, \textit{Science}, \textit{Blacksky}, \textit{AcademicSky}, entre outros), totalizando \textbf{168.463 postagens}. Em seguida, aplica \textit{Reservoir Sampling} sobre o arquivo de interações (\texttt{graphs/reposts.csv}), que contém mais de 152 milhões de linhas, selecionando uma amostra estatisticamente representativa sem carregar o arquivo completo em memória. O dataset resultante foi disponibilizado publicamente no repositório \textit{Hugging Face Datasets} para garantir a reprodutibilidade da pesquisa.

A integração contínua de dados complementares foi realizada utilizando os scripts desenvolvidos no módulo de coleta (\texttt{collect.py}), consumindo a API do Bluesky para mapear relações de engajamento em tempo real (\textit{follows}, \textit{reposts} e \textit{likes}).

\section{Modelagem e Treinamento da IA (Deep Learning)}

A inteligência central do sistema repousa sobre a biblioteca \texttt{PyTorch Geometric} (PyG), focada em aprendizado profundo sobre grafos.

\subsection{Vetorização de Texto (Embeddings)}

Antes da convolução geométrica, o texto puro das postagens precisa ser convertido em tensores matemáticos. Utilizou-se a biblioteca \texttt{sentence-transformers} com o modelo \texttt{paraphrase-multilingual-mpnet-base-v2} para extrair as características semânticas (\textit{features}) de cada nó da rede. Esse modelo produz vetores densos de \textbf{768 dimensões} com suporte nativo a português, inglês e espanhol --- cobrindo os idiomas predominantes nas postagens coletadas do Bluesky. Dessa forma, cada postagem no grafo carrega um vetor que representa o significado semântico do que foi escrito.

\subsection{Arquitetura da Rede Convolucional de Grafos}

A arquitetura da Rede Neural de Grafos implementada é composta por três camadas de convolução (\texttt{GCNConv}), projetadas para agregar informações de até três ``saltos'' (\textit{hops}) de distância na rede, capturando o contexto de amigos de amigos. Cada camada opera sobre representações de 64 dimensões no espaço latente. A regra de propagação camada a camada é:

\begin{equation}
H^{(l+1)} = \text{ReLU}\left(\tilde{D}^{-\frac{1}{2}} \tilde{A} \tilde{D}^{-\frac{1}{2}} H^{(l)} W^{(l)}\right)
\end{equation}

A camada final utiliza uma operação de \textit{Global Mean Pooling} para condensar todas as representações nodais do grafo em um único vetor de nível de grafo, seguida por uma camada \textit{Linear} com \textit{Dropout} de 50\% e classificação binária: \textit{Real} (0) ou \textit{Fake} (1).

\subsection{Pipeline de Treinamento e o Framework de Ensemble}

Para mitigar o problema do \textit{Cold Start} e maximizar a robustez do sistema, adotou-se uma estratégia de \textit{Ensemble} com quatro modelos GCN treinados em condições distintas, cada qual com um nível diferente de ceticismo sobre os dados do Bluesky:

\begin{itemize}
    \item \textbf{UPFD Baseline} (peso 10\%): treinado exclusivamente no \textit{benchmark} PolitiFact/Twitter. Representa o conhecimento histórico de redes maduras, mas com baixa capacidade de generalização para o Bluesky.
    \item \textbf{Bluesky Controlado} (peso 20\%): treinado com dados nativos do Bluesky em condições balanceadas. Oferece uma perspectiva intermediária entre o domínio do Twitter e a realidade da nova rede.
    \item \textbf{Bluesky Cético} (peso 30\%): treinado com penalidade sêxtupla (6x) sobre falsos positivos, tornando o modelo mais criterioso na classificação de \textit{fake news}.
    \item \textbf{Bluesky Exageradamente Cético} (peso 40\%): configuração de máxima penalidade. Recebe o maior peso no \textit{ensemble} por apresentar a menor taxa de falsos negativos nos experimentos --- ou seja, raramente classifica uma \textit{fake news} como real.
\end{itemize}

O veredito final é obtido pela soma ponderada das probabilidades de cada modelo:

\begin{equation}
S_{ensemble} = 0{,}10 \cdot P_{UPFD} + 0{,}20 \cdot P_{Ctrl} + 0{,}30 \cdot P_{C\acute{e}tico} + 0{,}40 \cdot P_{ExCetico}
\end{equation}

Se $S_{ensemble} \geq 0{,}5$, a publicação é classificada como \textit{Fake News}.

\section{Desenvolvimento da Aplicação Prática}

Para demonstrar a viabilidade da técnica em um cenário real, desenvolveu-se uma interface completa de interação.

\subsection{Backend Assíncrono e API}

A comunicação entre a Inteligência Artificial e a interface web é gerenciada por um servidor \textit{Backend} desenvolvido em \texttt{FastAPI}. Para evitar que operações computacionalmente intensas (geração de \textit{embeddings} BERT e inferência GNN) bloqueiem a interface do usuário, implementou-se um modelo de processamento assíncrono baseado em \textit{Background Tasks}:

\begin{enumerate}
    \item O cliente envia a URL de um post para o endpoint \texttt{POST /api/analyze} e recebe imediatamente um identificador único (\texttt{task\_id}), sem aguardar o processamento.
    \item O servidor executa a pipeline completa --- extração de texto via \textit{AT Protocol}, geração de \textit{embedding} BERT, construção do grafo e inferência do \textit{ensemble} --- em uma \textit{thread} separada do processo principal.
    \item O cliente realiza consultas periódicas ao endpoint \texttt{GET /api/result/\{task\_id\}} até que o status mude de \texttt{processing} para \texttt{done}.
\end{enumerate}

Os resultados de cada análise são persistidos em um banco de dados SQLite via \texttt{SQLAlchemy}, acessível por meio do endpoint \texttt{GET /api/history}, que exibe o histórico das dez consultas mais recentes.

\subsection{Frontend Web}

A interface de usuário foi concebida utilizando os \textit{frameworks} \texttt{React} e \texttt{Next.js} (versão 14). O portal web oferece uma experiência visual fluida onde o usuário pode colar o link de qualquer postagem pública do Bluesky e visualizar não apenas o veredito do \textit{ensemble}, mas também a contribuição individual de cada um dos quatro modelos, tornando transparente o processo de classificação preditiva baseado em grafos.
